\chapter{App G --- Activity Detail and Quality Assurance}

\section{App G --- Activity Detail and QA}

This appendix expands the phase summaries in Chapter with full activity lists, quality assurance procedures, and documentation and version control practices.

\subsection{Detailed Phase Activity Lists}

\subsubsection{Phase 0: Preparation and Infrastructure (July--November 2025)}

\begin{itemize}
\tightlist
\item
  HomeLab infrastructure initialisation (Proxmox, networking, storage)
\item
  Development environment setup (VSCode, LaTeX, Git)
\item
  Security tools installation (Parrot OS Security, Metasploit, Burp Suite)
\item
  Technical documentation and baseline establishment
\end{itemize}

\textbf{Deliverables}: Fully operational virtualisation infrastructure; development environment configuration; baseline documentation of system architecture.

\textbf{Duration}: 105 hours

\subsubsection{Phase 1: Preliminary Project and Conceptual Design (October 2025--January 2026)}

\begin{itemize}
\item
  \textbf{SMART Objectives Definition} (4h) --- Specification of Main Objective (MO): develop a reusable teaching package. Definition of 4 Specific Objectives (OE1--OE4) with measurable KPIs. Alignment with institutional GEISI curriculum requirements.
\item
  \textbf{State of the Art Analysis} (15h) --- Comprehensive review of 7 major cybersecurity training platforms (HackTheBox, TryHackMe, Cisco Academy, SEED Labs, GOAD, VulnHub, Hack4u). Identification of pedagogical gaps. Comparative analysis matrix across multiple dimensions.
\item
  \textbf{Functional and Non-Functional Requirements} (10h) --- RF1: design and implement 4 EVE-NG topologies (.unl files). RF2: implement vulnerabilities per PTES, OWASP, NIST, and CEH standards. RF3: create structured assessment materials and rubrics. NFR1--NFR3 covering performance, framework compliance, and institutional sustainability.
\item
  \textbf{Pentesting Methodology and Validation Framework} (8h) --- Adaptation of PTES for educational context. Definition of ethical hacking principles (EC-Council CEH). Mapping of attack phases to lab objectives. Validation criteria for security effectiveness.
\item
  \textbf{Design Decisions Justification} (5h) --- Rationale for EVE-NG selection vs alternative platforms. Technology stack justification (Parrot OS Security, Metasploit, Burp, OWASP ZAP). Pedagogical approach: hands-on labs with automation.
\item
  \textbf{Risk Analysis and Mitigation Planning} (15h) --- Identification of 12+ project risks (technical, pedagogical, institutional). Assessment of impact and probability. Definition of mitigation strategies.
\item
  \textbf{Resource Inventory and Allocation} (5h) --- Hardware requirements: dual-socket server with Proxmox. Software stack: EVE-NG Community Edition, Parrot OS Security images, security tools. Human resources: student ( 820h total), faculty supervisor. Budget considerations and open-source alternatives.
\item
  \textbf{Gantt Chart and Critical Path Analysis} (20h) --- Development of detailed project timeline with 37 tasks. Identification of 13 critical path tasks. CSV export and interactive web-based visualisation. Progress tracking and milestone definition.
\end{itemize}

\textbf{Deliverables}: Formal project proposal (Preliminary Project); detailed Gantt chart with critical paths; risk register and mitigation plan; requirements specification (functional and non-functional).

\textbf{Duration}: 82 hours

\subsubsection{Phase 2: Implementation and Interim Validation (January--March 2026)}

\begin{itemize}
\item
  \textbf{Lab 1: Reconnaissance Development} (14h research + 20h dev) --- Scenario: intelligence gathering phase using PTES framework. Tools: Nmap, passive reconnaissance techniques. Topology: multi-tier network with passive monitoring capabilities. Validation: successful execution of reconnaissance workflow.
\item
  \textbf{Lab 2: Web Vulnerabilities Development} (14h research + 25h dev) --- Scenario: OWASP Top 10 vulnerability exploitation. Tools: Burp Suite, OWASP ZAP, manual testing techniques. Topology: web application server with vulnerable services. Validation: exploitation success and reporting accuracy.
\item
  \textbf{Lab 3: Incident Response Development} (14h research + 25h dev) --- Scenario: log forensics and incident response (NIST SP 800-61 standards). Tools: Wireshark, tcpdump, network enumeration tools. Topology: multi-segment network with pre-staged attack artefacts. Validation: successful timeline reconstruction and flag recovery.
\item
  \textbf{Pilot Validation with Users} (20h + 6h revision) --- Informal testing sessions with 4--5 GEISI students conducted as each lab is completed. Usability observations and timing notes. Post-session survey on satisfaction, difficulty, and guide clarity. Adjustments applied based on participant feedback.
\item
  \textbf{Topology Adjustments and Optimisation} (30h) --- Performance tuning based on pilot feedback. VM image optimisation and deployment time reduction. Network configuration refinement for stability.
\item
  \textbf{Teaching Material Production} (40h) --- Student lab guide PDF (student lab guide) per lab: task description, objectives, hints, and tool reference. Instructor resolution guide PDF (instructor solution guide) per lab: one possible walkthrough, flag values, assessment rubric, and common errors. Technical configuration reference per lab in the repository.
\end{itemize}

\textbf{Deliverables}: 3 fully functional EVE-NG labs; student and instructor guides; pilot validation report; technical configuration references.

\textbf{Duration}: 188 hours

\subsubsection{Phase 3: Completion and Final Validation (March--May 2026)}

\begin{itemize}
\item
  \textbf{Lab 4: CipherStrike --- Cryptography \& Steganography} (16h research + 30h dev) --- Scenario: multi-server CTF with cryptography (ROT13, AES-ECB, Vigen\'{e}re, XOR, RSA small exponent) and steganography (EXIF, LSB, steghide) challenges across three servers with a gate mechanism. Tools: Python cryptanalysis, steghide, exiftool, gmpy2. Topology: ServerA (crypto), ServerB (stego), ServerC (gated advanced mix), Defender (SOC monitoring). Validation: 14 flags verified end-to-end, gate mechanism (MD5 of A5+B5) tested.
\item
  \textbf{Comprehensive Penetration Testing} (30h Round 1 + 25h Round 2) --- Round 1: sequential testing of Labs 1, 2, and 3. Round 2: integrated testing across all 4 laboratories. Documentation of all findings and remediation strategies.
\item
  \textbf{Automation Scripts Development} (30h) --- Utility scripts (Bash/Python) for recurring lab management tasks: bridge configuration, ISO upload, and connectivity check scripts. Scripts documented and tested in the project EVE-NG environment.
\item
  \textbf{Final User Validation and Iteration} (18h) --- Final informal validation sessions with the pilot group (4--5 participants). Collection of post-session survey data and usability notes. Implementation of final adjustments based on feedback.
\item
  \textbf{Bachelor's Thesis Report Writing and Finalisation} (90h) --- Phase I (15h): Introduction, Objectives, Literature Review. Phase II (40h): Methodology, Requirements, Implementation Results. Phase III (35h): Validation, Analysis, Conclusions. Final revision, bibliography integration, and publication preparation (35h).
\end{itemize}

\textbf{Deliverables}: 4 fully functional labs; utility automation scripts; final bachelor's thesis report; pilot validation report; deployment package ready for institutional use.

\textbf{Duration}: 315 hours

\subsection{Quality Assurance Procedures}

\subsubsection{Technical Validation}

\begin{itemize}
\tightlist
\item
  Automated testing of lab deployment and reset
\item
  Manual security testing and penetration verification
\item
  Performance benchmarking against KPI targets
\item
  Compatibility testing across EVE-NG versions
\end{itemize}

\subsubsection{Pedagogical Validation}

\begin{itemize}
\tightlist
\item
  Alignment with GEISI learning outcomes
\item
  Assessment of student understanding through rubrics
\item
  Effectiveness in achieving educational objectives
\item
  Feedback from faculty and student participants
\end{itemize}

\subsubsection{Usability Testing}

\begin{itemize}
\tightlist
\item
  User interface and documentation clarity
\item
  Accessibility and ease of deployment
\item
  Support and documentation adequacy
\item
  User satisfaction scoring (≥ 4/5 target)
\end{itemize}

\subsection{Documentation and Version Control}

\begin{itemize}
\tightlist
\item
  GitHub repository (\texttt{TheEgea/TFG}) for all code, scripts, and documentation, licensed under CC BY-NC-SA 4.0
\item
  LaTeX (XeLaTeX + OpenDyslexic) for formal academic documentation (Memory and Appendices)
\item
  MkDocs Material for the operational web documentation site
\item
  Versioning system for all deliverables with semantic commit messages
\item
  Change logs and improvement tracking through pull requests
\item
  Deployment guides and runbooks per lab
\end{itemize}
